\section{如何选择合适的模式}

每种模式回答的是不同的问题。以下是选择指南：

\begin{table}[H]
\centering
\begin{tabular}{@{}lll@{}}
\toprule
\textbf{模式} & \textbf{核心问题} & \textbf{典型场景} \\
\midrule
Tool Wrapper & 如何注入特定库的知识？ & 编码规范、框架约定 \\
Generator & 如何保证输出结构一致？ & API 文档、项目脚手架 \\
Reviewer & 如何系统化检查代码？ & PR review、安全审计 \\
Inversion & 如何确保信息收集完整？ & 项目规划、架构设计 \\
Pipeline & 如何强制多步骤顺序执行？ & 文档流水线、审批工作流 \\
\bottomrule
\end{tabular}
\caption{五种 Agent Skill 设计模式对比}
\end{table}

\begin{knowledgebox}{模式可以组合}
这五种模式并非互斥。它们可以自由组合：
\begin{itemize}
    \item Pipeline 可以在最后一步嵌入 Reviewer 来自检工作成果
    \item Generator 可以在开始阶段使用 Inversion 来收集必要的模板变量
\end{itemize}
得益于 ADK 的 \texttt{SkillToolset} 和渐进式披露（progressive disclosure）机制，Agent 在运行时只在需要的模式上花费 context token。
\end{knowledgebox}

\subsection{本章小结}

选择模式时应根据核心问题匹配，而非盲目套用。更重要的是，模式之间可以组合使用，形成更强大的复合 skill。

